% Auto-generated by paper/make_v3_tables.py -- do not hand-edit.
\begin{table}[t]
\centering
\small
\begin{tabular}{llrrr}
\toprule
Arm & Task & $\sum$ parts & \texttt{all} & ratio \\
\midrule
\multicolumn{5}{l}{\emph{other tasks}} \\
1.7b & connected\_nodes & $+5.7$ & $+2.0$ & $0.35$ \\
1.7b & edge\_existence & $+12.8$ & $+11.5$ & $0.90$ \\
1.7b & node\_degree & $+14.5$ & $+10.2$ & $0.71$ \\
1.7b-t & connected\_nodes & $+8.8$ & $+6.2$ & $0.71$ \\
1.7b-t & edge\_existence & $-4.5$ & $-2.0$ & $0.44$ \\
1.7b-t & node\_degree & $+16.5$ & $+12.5$ & $0.76$ \\
4b & edge\_existence & $-10.8$ & $-10.0$ & $0.93$ \\
4b & node\_degree & $-9.2$ & $-8.0$ & $0.86$ \\
\midrule
\multicolumn{5}{l}{\emph{\texttt{edge\_count}}} \\
1.7b-t & edge\_count & $-8.0$ & $+9.0$ & $-1.12$ \\
4b & edge\_count & $+24.5$ & $+3.7$ & $0.15$ \\
4b-t & edge\_count & $+35.2$ & $-6.7$ & $-0.19$ \\
\midrule
\multicolumn{4}{l}{\emph{median over cells}} & $0.44$ \\
\bottomrule
\end{tabular}
\caption{\texttt{all} renders \texttt{degree}, \texttt{clustering} and \texttt{rwse} in one primer, so the parts and the bundle are both measured. The first two columns are paired differences against \texttt{none} in percentage points, averaged over $p \in \{0.10, 0.20, 0.35, 0.50\}$; rows whose parts sum to under $4$ points are omitted, since a ratio of two null effects is not informative. The final row is the median over the $48$ individual (arm, task, density) cells rather than over the averaged rows above, because averaging numerator and denominator over densities first gives a ratio of means that the largest cells dominate; weighted the other ways the estimate runs from $0.29$ to $0.71$, so it bounds a tendency and is not a coefficient. The bootstrap interval on the per-cell median is $[0.33, 0.63]$, which excludes $1$: the bundle delivers reliably less than the sum of its parts. Against the \emph{largest} single part instead, the median is $0.80$ $[0.56, 1.08]$, which does not exclude $1$ (Section~\ref{sec:results-composition}). \texttt{edge\_count} is the extreme case: there the bundle keeps almost nothing and in two arms moves against its parts.}
\label{tab:additivity}
\end{table}
